\documentclass[11pt,a4paper]{article}
\usepackage[margin=23mm]{geometry}
\usepackage{amsmath,amssymb,amsthm}
\setlength{\emergencystretch}{2em}
\newtheorem{theorem}{Theorem}
\newcommand{\gen}[1]{\langle#1\rangle}
\newcommand{\Alt}{\operatorname{Alt}}
\title{A four-point internal orbit gives a nonterminal separator}
\author{Human Compression project: L044 continuation}
\date{}
\begin{document}\maketitle
Let $G=\gen{\rho_0,\ldots,\rho_{r-1}}\le\Alt(\Omega)$
be transitive with a connected string C-marking and all
maximal omissions intransitive. Assume $n=2r+1$ and the
actual six-point head on $D=\{1,\ldots,6\}$, with
\[
 \tau=\rho_3=(6\ b)\delta,\qquad
 R=G_{\ge3}=\Alt(\Omega\setminus\{1,\ldots,5\}).
\]
The involution $\delta$ fixes $D,b$; $\sigma=\rho_4$
fixes $D$; $H=G_{\ge5}$ fixes $D,b$.
Put $K=G_{>5}$. The actual pure-overlap argument gives
$x=b^\sigma\ne b$, $\delta x=x$ and $Kx=x$.

Suppose $H$ has precisely the seven fixed points $D\cup\{b\}$
and three nontrivial orbits $E,W,V$. Retain the following
literal assignment data:
\begin{itemize}
\item $|E|=2$, $|W|=4$, and $|V|$ is even. The involution
      $\delta$ swaps $E$, is free on $W$, and fixes $V$.
\item The normalized assigned sets are $F_E=\varnothing$,
      $F_W=\{j\}$ for some $j\ge6$, and
      $F_V=\{5,\ldots,r-1\}\setminus\{j\}$.
      Each assigned colour has two full ambient-eligible
      edges in its assigned orbit.
\item $K$ is intransitive on $W$.
\end{itemize}

\begin{theorem}
The index $j$ is strictly below $r-1$, and
\[
                       |\operatorname{Fix}(G_{>j})|\ge2j+2.
\]
Thus this precise four-point-orbit budget has a nonterminal
local separator. No assertion about its transfer through a
subsequent external attachment is made.
\end{theorem}
\begin{proof}
The two pairs of $\delta|_W$ are preserved by $H$.
An eligible $j$-edge cannot be one of these pairs,
because $\tau\in G_j$ joins that pair. Thus $\rho_j$
exchanges the pair vertices. Every other mark of $H$
has trivial pair image: if one also exchanged them,
it and $\tau$ would connect all four points inside
$G_j$, contradicting the eligible edge.

The group $K$ contains $\rho_j$, so it is transitive
on the two pair vertices. As it is intransitive on $W$,
it has exactly two orbits of size two, each meeting
each pair once. A pair-trivial element of $K$ fixes
both orbits and hence all four points. Consequently
$K|_W$ is the diagonal $S_2$ on these two rows,
$\rho_j|_W$ is a double transposition, and every
other mark of $K$ fixes $W$ pointwise. This proves
the needed two-pair equality directly; no theorem
stated only for at least three pairs is imported.

Write $Q=\gen{\rho_j}$ and $Z_\ell=W^{G_{\le\ell}}$
for $4\le\ell\le j-2$, where the superscript denotes
the union of all translates. The prefix $G_{\le4}$
centralizes $K$, so $Z_4$ is a union of its two-point
row orbits. It contains at least three such rows.
Indeed, if $Z_4=W$, both $G_{\le4}$ and $H$ would
preserve $W$, contradicting transitivity of $G$.

Every $Z_\ell$ is a union of nontrivial two-point
$Q$-orbits: $G_{\le\ell}$ centralizes $Q$ in the
stated range. These sets exclude the $Q$-fixed head
and are therefore proper. For $4\le\ell\le j-3$,
each mark of index at least $\ell+2$, other than
$\rho_j$, fixes $W$ and commutes with $G_{\le\ell}$,
so fixes $Z_\ell$ pointwise. The mark $\rho_j$
preserves it, as does $G_{\le\ell}$. Thus
$\rho_{\ell+1}$ cannot preserve $Z_\ell$, or the
whole group would preserve a proper set. It adds at
least one complete two-point $Q$-orbit. Hence
\[
                       |Z_{j-2}|\ge6+2(j-6)=2j-6.       \tag{1}
\]
Every mark after $j$ fixes $W$ and commutes with
$G_{\le j-2}$, so fixes $Z_{j-2}$. It also fixes
the eight distinct points $D\cup\{b,x\}$, since
it belongs to $K$. These points are fixed by $Q$
and therefore disjoint from $Z_{j-2}$. Adding them
to (1) proves the required bound.

Suppose now that $j=r-1$, so $n=2j+3$.
Let $C_k$ be the actual $G_{\le k}$-orbit containing
$D$. The degree-independent support-merging proof
of the acquired prefix-kernel theorem gives
\[
 \Alt(C_k)\le G_{\le k},\qquad |C_k|\ge k+4\quad(k\ge3).
\]
Here is the precise part used. The literal head gives
pure $\Alt(D\cup\{b\})$ and then
pure $\Alt(D\cup\{b,x\})$. Each subsequent mark
fixes $C_{k-2}$ by string commutation and its pure
alternating kernel. Merging the overlapping conjugate
supports supplies pure $\Alt(C_k)$. If a head orbit
failed to grow, every later mark would fix it by the
same commutation, contradicting transitivity or
independence. These arguments use neither $n=2r$
nor an alternating rank induction.

Since $Q$ centralizes $G_{\le j-2}$, it fixes
$C_{j-2}$ pointwise. This set is disjoint from
$Z_{j-2}$, all of whose points are moved by $Q$.
Consequently
\[
                  2j+3=n\ge(2j-6)+(j+2)=3j-4,
\]
so $j\le7$. There are only two possibilities.
If $j=6$, then $r=7$, $n=15$ and
$|V|=n-7-2-4=2$. But $F_V=\{5\}$ requires two
disjoint eligible edges on this two-point set,
which is impossible. If $j=7$, then $r=8$,
$n=17$, $|V|=4$ and $F_V=\{5,6\}$.
The two eligible $6$-edges on $V$ make
$\rho_6|_V$ fixed-point-free. On the other hand
$x\in V$, because it is outside $D,b$ and is
fixed by $\delta$, which is free on $E\cup W$.
The group $K$ fixes $x$, and contains $\rho_6$.
This is a contradiction. Thus $j<r-1$.
\end{proof}

\paragraph{Attribution and boundary.}
The row-transport and prefix-growth strategy is acquired;
the two-pair equality, its bound $2j+2$, and the terminal
labels six and seven are proved here. The pure alternating
kernel is the acquired support-merging argument, retaining
the root collaborator's initial kernel contribution and
all complete predecessor costs. The three-orbit assignment
budget is explicit input; neither its necessity in every
seven-fixed-point configuration nor global separator
transfer is asserted. No new computation, formal theorem,
original or question credit is claimed.
\end{document}
